\section{Graph-Based Multi-Date Crown Tracking}

This section presents the core algorithmic framework for tracking individual tree crowns across multiple orthomosaics. We formulate crown tracking as a graph matching problem, where nodes represent crown detections at specific time points and edges encode temporal correspondences. This graph-based approach naturally handles complex scenarios including crown splits, merges, appearances, disappearances, and ambiguous matches that arise from detection inconsistencies and phenological changes.

%--------------------------------------------------------------------
\subsection{High-Level Pipeline Overview}

The crown tracking pipeline operates on a sequence of $N$ orthomosaics $\{\mathcal{O}_1, \mathcal{O}_2, \ldots, \mathcal{O}_N\}$ captured at discrete time points $t_1 < t_2 < \cdots < t_N$. For each orthomosaic $\mathcal{O}_i$, we obtain a set of crown polygons $\mathcal{C}_i = \{c_{i,j}\}_{j=1}^{m_i}$ via the Detectree2 segmentation pipeline described in Section~3. The tracking problem seeks to establish correspondences between crowns across consecutive time points, ultimately constructing temporal chains that represent individual trees tracked through time.

\paragraph{Pipeline stages:}
\begin{enumerate}[leftmargin=*, itemsep=2pt]
    \item \textbf{Spatial alignment} (Section~3): Consecutive orthomosaics are aligned using translational shifts computed from overlapping crown centroids to minimize georeferencing drift.
    
    \item \textbf{Attribute computation}: For each crown $c_{i,j} \in \mathcal{C}_i$, extract geometric attributes:
    \begin{itemize}
        \item Centroid coordinates $(x_{i,j}, y_{i,j})$
        \item Area $A_{i,j}$ (m$^2$) and perimeter $P_{i,j}$ (m)
        \item Compactness $\kappa_{i,j} = 4\pi A_{i,j} / P_{i,j}^2 \in [0,1]$
        \item Eccentricity $\epsilon_{i,j}$ from minimum rotated rectangle
        \item Bounding box and aspect ratio
    \end{itemize}
    
    \item \textbf{Pairwise candidate generation}: For each consecutive pair $(\mathcal{O}_i, \mathcal{O}_{i+1})$, identify all crown pairs $(c_{i,j}, c_{i+1,k})$ within a spatial proximity threshold $d_{\text{max}}$ (typically 75--200~m).
    
    \item \textbf{Feature computation}: For each candidate pair, compute similarity metrics (Section~\ref{sec:pairwise_features}) including IoU, overlap ratios, centroid distance, area similarity, and shape similarity.
    
    \item \textbf{Case classification}: Assign each candidate to one of four match cases—\texttt{one\_to\_one}, \texttt{containment}, \texttt{nearby}, or \texttt{none}—based on conditional thresholds (Section~\ref{sec:case_classification}).
    
    \item \textbf{Edge selection}: For each case in priority order, select the best-scoring edges subject to cardinality constraints (at most one outgoing edge per previous crown, one incoming edge per current crown for strict matching).
    
    \item \textbf{Graph construction}: Build a directed graph $G = (V, E)$ where:
    \begin{itemize}
        \item $V = \bigcup_{i=1}^N \{(i,j) : c_{i,j} \in \mathcal{C}_i\}$ (nodes indexed by $(om\_id, crown\_index)$)
        \item $E \subseteq V \times V$ with edges $(u, v)$ only if $u = (i,j)$ and $v = (i+1,k)$ (forward temporal flow)
    \end{itemize}
    
    \item \textbf{Chain extraction}: Extract all maximal paths (chains) from source nodes (in-degree 0) to sink nodes (out-degree 0), representing complete tracking histories.
    
    \item \textbf{Quality assessment}: Compute match rates, chain length distributions, and graph complexity metrics to evaluate tracking performance.
\end{enumerate}

The entire pipeline is implemented in the \texttt{TreeTrackingGraph} Python class, which provides a modular interface for configuring matching parameters, building graphs with different presets (strict vs.\ relaxed), and exporting quality reports.

%--------------------------------------------------------------------
\subsection{Limitations of Hungarian Matching}

Before introducing our graph-based approach, we motivate its necessity by examining the limitations of classical bipartite matching algorithms, particularly the Hungarian algorithm~\cite{kuhn1955hungarian}, which has been widely used for object tracking in computer vision.

%--------------------------------------------------------------------
\subsubsection{One-to-One Constraint and Split/Merge Failures}

The Hungarian algorithm solves the \emph{linear assignment problem}: given a bipartite graph with crowns from $\mathcal{C}_i$ on one side and crowns from $\mathcal{C}_{i+1}$ on the other, find a maximum-weight perfect matching. This enforces a strict one-to-one correspondence: each crown in $\mathcal{C}_i$ matches to at most one crown in $\mathcal{C}_{i+1}$, and vice versa.

\paragraph{Crown splits:} When a single crown $c_{i,j}$ in $\mathcal{O}_i$ is over-segmented into two crowns $c_{i+1,k}$ and $c_{i+1,\ell}$ in $\mathcal{O}_{i+1}$ (due to improved resolution, changed viewpoint, or segmentation artifacts), both $c_{i+1,k}$ and $c_{i+1,\ell}$ may have high overlap with $c_{i,j}$. The Hungarian algorithm is forced to select only one match, discarding the other and breaking the tracking chain.

\paragraph{Crown merges:} Conversely, when two adjacent crowns $c_{i,j}$ and $c_{i,k}$ under-segment into a single crown $c_{i+1,\ell}$, both may have substantial overlap with $c_{i+1,\ell}$. Again, the one-to-one constraint prevents representing this many-to-one relationship.

\paragraph{Empirical impact:} In our initial experiments with Hungarian matching on the SIT dataset, we observed frequent chain breaks at split/merge events, with effective chain lengths reduced by 30--50\% compared to the graph-based approach that explicitly models these cases.

%--------------------------------------------------------------------
\subsubsection{Local Pairwise Optimization Issues}

The Hungarian algorithm optimizes a global objective over a single bipartite graph:
\[
\max_{\pi} \sum_{j} w(c_{i,j}, c_{i+1,\pi(j)})
\]
subject to $\pi : \mathcal{C}_i \to \mathcal{C}_{i+1}$,
where $w(\cdot,\cdot)$ is a similarity weight (e.g., IoU). However, this optimization is \emph{local} to the pair $(\mathcal{O}_i, \mathcal{O}_{i+1})$ and ignores information from other time points.

\paragraph{Transitive inconsistencies:} Consider three consecutive orthomosaics $\mathcal{O}_i, \mathcal{O}_{i+1}, \mathcal{O}_{i+2}$. The Hungarian algorithm independently computes:
\begin{itemize}
    \item $\pi_{i,i+1}: \mathcal{C}_i \to \mathcal{C}_{i+1}$ (matches for $\mathcal{O}_i \to \mathcal{O}_{i+1}$)
    \item $\pi_{i+1,i+2}: \mathcal{C}_{i+1} \to \mathcal{C}_{i+2}$ (matches for $\mathcal{O}_{i+1} \to \mathcal{O}_{i+2}$)
\end{itemize}
There is no guarantee that the composed mapping $\pi_{i+1,i+2} \circ \pi_{i,i+1}$ is consistent. For example, if $c_{i,j} \to c_{i+1,k}$ and $c_{i+1,k} \to c_{i+2,\ell}$ are both selected, but $c_{i,j}$ has no direct correspondence with $c_{i+2,\ell}$ (e.g., due to intermediate detection failure), the chain $c_{i,j} \to c_{i+1,k} \to c_{i+2,\ell}$ may not represent the same physical tree.

\paragraph{Myopic feature selection:} Pairwise matching cannot leverage temporal context. A weak match between $\mathcal{O}_i$ and $\mathcal{O}_{i+1}$ might be disambiguated by examining the crown's trajectory across $\mathcal{O}_{i-1}, \mathcal{O}_i, \mathcal{O}_{i+1}$, but the Hungarian framework provides no mechanism to incorporate such multi-date consistency checks.

%--------------------------------------------------------------------
\subsubsection{Inconsistent Chains Across Multiple Dates}

When applying Hungarian matching sequentially across $N-1$ pairs of consecutive orthomosaics, the resulting chains often exhibit logical inconsistencies:
\begin{itemize}
    \item \textbf{Branching chains}: A crown $c_{i,j}$ matches to $c_{i+1,k}$, which then matches to both $c_{i+2,\ell}$ and $c_{i+2,m}$ (if the algorithm is run separately for $\mathcal{O}_{i+1} \to \mathcal{O}_{i+2}$, the one-to-one constraint is applied per pair, not globally).
    
    \item \textbf{Broken chains}: A tree tracked as $c_{1,j} \to c_{2,k} \to c_{3,\ell}$ may suddenly match to a different crown $c_{4,m}$ at the fourth time point due to a locally optimal but globally inconsistent choice.
    
    \item \textbf{Identity switches}: Two crowns may swap identities across intermediate time points, a phenomenon well-documented in multi-object tracking literature~\cite{leal2016motchallenge}.
\end{itemize}

These issues motivate a holistic graph-based approach where all candidate edges across all pairs are represented simultaneously, and selection is performed with global awareness of chain structure.

%--------------------------------------------------------------------
\subsection{From Pairwise Matching to Crown Graphs}

Our graph-based tracking framework addresses the limitations of Hungarian matching by constructing a unified directed graph $G = (V, E)$ that encodes all plausible temporal correspondences.

\paragraph{Node set:} Each node $v \in V$ uniquely identifies a crown detection at a specific time point:
\[
V = \{(i, j) : i \in \{1, 2, \ldots, N\}, \ c_{i,j} \in \mathcal{C}_i\}
\]
where $(i,j)$ denotes the $j$-th crown in the $i$-th orthomosaic. Node attributes include all geometric properties of $c_{i,j}$ (centroid, area, compactness, etc.).

\paragraph{Edge set:} An edge $(u, v) \in E$ represents a candidate temporal correspondence between crowns. We restrict edges to \emph{forward consecutive pairs}:
\[
E \subseteq \{((i,j), (i+1,k)) : i \in \{1, \ldots, N-1\}, \ c_{i,j} \in \mathcal{C}_i, \ c_{i+1,k} \in \mathcal{C}_{i+1}\}
\]
This enforces temporal causality (no backward edges) and limits complexity by considering only adjacent time points. Each edge carries attributes:
\begin{itemize}
    \item \textbf{Similarity score} $s(u,v) \in [0,1]$: weighted combination of geometric and overlap features (Section~\ref{sec:pairwise_features}).
    \item \textbf{Case label} $\text{case}(u,v) \in \{\texttt{one\_to\_one}, \texttt{containment}, \texttt{nearby}, \texttt{none}\}$: match type (Section~\ref{sec:case_classification}).
    \item \textbf{Overlap metrics}: IoU, overlap ratios, centroid distance.
    \item \textbf{Component similarities}: spatial, area, shape, and IoU components of $s(u,v)$.
\end{itemize}

\paragraph{Graph construction vs.\ Hungarian:} Unlike the Hungarian approach, which solves $N-1$ independent bipartite matching problems, we construct $G$ by:
\begin{enumerate}
    \item Generating all candidate edges $(u,v)$ with $s(u,v)$ above a threshold.
    \item Classifying edges into match cases based on overlap patterns.
    \item Selecting edges \emph{globally} across all pairs using a case-prioritized greedy algorithm with cardinality constraints.
\end{enumerate}
This allows us to:
\begin{itemize}
    \item Model splits/merges by permitting nodes with multiple incoming or outgoing edges (relaxed in ultra-relaxed preset).
    \item Incorporate case-specific selection rules (e.g., mutual best matching for high-confidence \texttt{one\_to\_one} edges).
    \item Analyze chain structure (paths through $G$) to validate consistency.
\end{itemize}

\paragraph{Chains and paths:} A \emph{tracking chain} is a directed path $\mathcal{P} = (v_1, v_2, \ldots, v_\ell)$ in $G$ where:
\begin{itemize}
    \item $v_1 = (i_1, j_1)$ is a source (in-degree 0).
    \item $v_\ell = (i_\ell, j_\ell)$ is a sink (out-degree 0).
    \item $(v_k, v_{k+1}) \in E$ for all $k \in \{1, \ldots, \ell-1\}$.
    \item $i_1 < i_2 < \cdots < i_\ell$ (monotonically increasing time indices).
\end{itemize}
The length $\ell$ of the chain indicates how many orthomosaics the tree was successfully tracked through. A \emph{full-length chain} has $\ell = N$ (present in all orthomosaics).

%--------------------------------------------------------------------
\subsection{Pairwise Matching Features}
\label{sec:pairwise_features}

For each candidate crown pair $(c_{i,j}, c_{i+1,k})$, we compute a rich feature vector that quantifies spatial, geometric, and topological similarity. These features form the basis for case classification and edge scoring.

%--------------------------------------------------------------------
\subsubsection{Overlap-Based Features: IoU and Overlap Ratios}

Let $A_{i,j}$ and $A_{i+1,k}$ denote the areas of crowns $c_{i,j}$ and $c_{i+1,k}$, respectively, and let $I_{jk} = c_{i,j} \cap c_{i+1,k}$ be their spatial intersection.

\paragraph{Intersection over Union (IoU):}
\[
\text{IoU}(c_{i,j}, c_{i+1,k}) = \frac{|I_{jk}|}{|c_{i,j} \cup c_{i+1,k}|} = \frac{|I_{jk}|}{A_{i,j} + A_{i+1,k} - |I_{jk}|}
\]
IoU is a symmetric metric in $[0,1]$, with $\text{IoU} = 1$ for perfect overlap and $\text{IoU} = 0$ for disjoint crowns. It is sensitive to both position and size: even crowns with identical shapes will have low IoU if significantly displaced.

\paragraph{Asymmetric overlap ratios:}
To capture directional containment relationships, we define:
\[
\text{overlap}_{\text{prev}}(c_{i,j}, c_{i+1,k}) = \frac{|I_{jk}|}{A_{i,j}}, \quad
\text{overlap}_{\text{curr}}(c_{i,j}, c_{i+1,k}) = \frac{|I_{jk}|}{A_{i+1,k}}
\]
These ratios measure the fraction of each crown covered by the intersection:
\begin{itemize}
    \item $\text{overlap}_{\text{prev}} \approx 1$ and $\text{overlap}_{\text{curr}} \ll 1$ suggests $c_{i+1,k}$ is a small sub-crown of $c_{i,j}$ (potential split).
    \item $\text{overlap}_{\text{prev}} \ll 1$ and $\text{overlap}_{\text{curr}} \approx 1$ suggests $c_{i,j}$ is subsumed by $c_{i+1,k}$ (potential merge).
    \item Both $\approx 1$ indicates strong one-to-one correspondence.
\end{itemize}

\paragraph{Overlap gate:} We introduce a global parameter $\tau_{\text{overlap}} \in [0,1]$ (typically 0.05--0.48 depending on preset) such that a crown pair is considered \emph{overlapping} only if:
\[
\max(\text{overlap}_{\text{prev}}, \text{overlap}_{\text{curr}}) \geq \tau_{\text{overlap}}
\]
This gate filters out weak spatial associations and reduces the candidate set.

%--------------------------------------------------------------------
\subsubsection{Geometric Features: Centroid Distance and Area Ratios}

\paragraph{Centroid distance:} Let $\mathbf{p}_{i,j} = (x_{i,j}, y_{i,j})$ and $\mathbf{p}_{i+1,k} = (x_{i+1,k}, y_{i+1,k})$ denote the centroids of $c_{i,j}$ and $c_{i+1,k}$. The Euclidean distance is:
\[
d_{\text{centroid}}(c_{i,j}, c_{i+1,k}) = \|\mathbf{p}_{i,j} - \mathbf{p}_{i+1,k}\|_2
\]
measured in meters (assuming projected CRS). We impose a hard constraint $d_{\text{centroid}} \leq d_{\max}$ to prune distant pairs. For pairs passing this filter, we compute a normalized spatial similarity:
\[
s_{\text{spatial}}(c_{i,j}, c_{i+1,k}) = \max\left(0, 1 - \frac{d_{\text{centroid}}}{d_{\max}}\right)
\]
This similarity decreases linearly with distance, reaching 0 at $d_{\max}$.

\paragraph{Area similarity:} Let $A_{i,j}$ and $A_{i+1,k}$ be the crown areas. We define:
\[
s_{\text{area}}(c_{i,j}, c_{i+1,k}) = \frac{\min(A_{i,j}, A_{i+1,k})}{\max(A_{i,j}, A_{i+1,k})} \in [0, 1]
\]
This symmetric similarity is 1 for identical areas and decreases as the size ratio diverges. It captures crown growth or shrinkage across time points.

\paragraph{Area ratio (asymmetric):} For case classification, we also compute the directional ratio:
\[
r_{\text{area}}(c_{i,j}, c_{i+1,k}) = \frac{A_{i+1,k}}{A_{i,j}}
\]
Ratios $r \gg 1$ suggest crown expansion or merging; $r \ll 1$ suggest shrinkage or splitting.

\paragraph{Mean radius:} To contextualize centroid distances relative to crown size, we compute an effective radius for each crown by assuming a circular approximation:
\[
R_{i,j} = \sqrt{\frac{A_{i,j}}{\pi}}
\]
The mean radius for the pair is:
\[
\bar{R}_{jk} = \frac{R_{i,j} + R_{i+1,k}}{2}
\]
A normalized centroid factor is then:
\[
f_{\text{centroid}}(c_{i,j}, c_{i+1,k}) = 1 - \min\left(1, \frac{d_{\text{centroid}}}{3\bar{R}_{jk}}\right)
\]
This penalizes distances beyond 3 mean radii (approximately 3 crown diameters).

%--------------------------------------------------------------------
\subsubsection{Containment Flags and Proximity Windows}

\paragraph{Geometric containment:} We compute boolean flags:
\begin{align*}
\text{prev\_contains\_curr} &= \mathbb{1}[c_{i+1,k} \subseteq c_{i,j}] \\
\text{curr\_contains\_prev} &= \mathbb{1}[c_{i,j} \subseteq c_{i+1,k}]
\end{align*}
using Shapely's \texttt{contains()} predicate (with buffering by $\epsilon = 0.01$~m to handle numerical precision). Containment indicates a significant change in crown boundary, often due to segmentation artifacts or real canopy growth/decline.

\paragraph{Shape similarity:} We compute two shape descriptors:
\begin{itemize}
    \item \textbf{Compactness similarity}: 
    \[
    s_{\text{compact}}(c_{i,j}, c_{i+1,k}) = 1 - |\kappa_{i,j} - \kappa_{i+1,k}|
    \]
    where $\kappa = 4\pi A / P^2$ (Polsby--Popper compactness).
    
    \item \textbf{Eccentricity similarity}:
    \[
    s_{\text{eccentric}}(c_{i,j}, c_{i+1,k}) = 1 - |\epsilon_{i,j} - \epsilon_{i+1,k}|
    \]
    where $\epsilon$ is the eccentricity (minor/major axis ratio) from the minimum rotated rectangle.
\end{itemize}
The overall shape similarity is:
\[
s_{\text{shape}}(c_{i,j}, c_{i+1,k}) = \frac{s_{\text{compact}} + s_{\text{eccentric}}}{2}
\]

\paragraph{Proximity window:} For the \texttt{nearby} case (ultra-relaxed preset), we define a tight proximity threshold $d_{\text{near}} = 30$~m. Crown pairs with $d_{\text{centroid}} < d_{\text{near}}$ are eligible for proximity-only matching even if $\text{IoU} \approx 0$ (accounting for large shape changes or segmentation drift).

%--------------------------------------------------------------------
\subsection{Conditional Threshold Matching}
\label{sec:conditional_matching}

Rather than using a single global similarity threshold, we adopt a \emph{conditional threshold} strategy where different match types are subject to tailored criteria. This is motivated by the observation that high-confidence one-to-one matches exhibit fundamentally different feature distributions than ambiguous merges or proximity-only fallbacks.

%--------------------------------------------------------------------
\subsubsection{Base Similarity and Overlap Gates}

For each crown pair $(c_{i,j}, c_{i+1,k})$, we first compute a \emph{base similarity} as a weighted combination of the component similarities:
\[
s_{\text{base}}(c_{i,j}, c_{i+1,k}) = \sum_{\alpha \in \{\text{spatial, area, shape, iou}\}} w_\alpha^{\text{case}} \cdot s_\alpha(c_{i,j}, c_{i+1,k})
\]
where the weights $w_\alpha^{\text{case}}$ depend on the match case (Section~\ref{sec:case_classification}). For example, \texttt{one\_to\_one} weights might be:
\[
w_{\text{spatial}}^{\text{1-1}} = 0.45, \quad w_{\text{area}}^{\text{1-1}} = 0.20, \quad w_{\text{shape}}^{\text{1-1}} = 0.15, \quad w_{\text{iou}}^{\text{1-1}} = 0.20
\]
emphasizing spatial alignment and IoU. In contrast, \texttt{nearby} weights are:
\[
w_{\text{spatial}}^{\text{near}} = 0.85, \quad w_{\text{area}}^{\text{near}} = 0.10, \quad w_{\text{shape}}^{\text{near}} = 0.03, \quad w_{\text{iou}}^{\text{near}} = 0.02
\]
heavily prioritizing spatial proximity over shape or overlap.

\paragraph{Early filtering:} Before case classification, we apply two global filters:
\begin{enumerate}
    \item \textbf{Spatial filter}: Reject pairs with $d_{\text{centroid}} > d_{\max}$ (typically 75--200~m).
    \item \textbf{Minimum quality filter}: Reject pairs where both $s_{\text{base}} < \tau_{\text{sim}}$ \emph{and} $\text{IoU} < \tau_{\text{overlap}}$, ensuring at least one of base similarity or overlap is non-negligible.
\end{enumerate}
These filters reduce the candidate set from $O(|\mathcal{C}_i| \times |\mathcal{C}_{i+1}|)$ to a manageable size.

%--------------------------------------------------------------------
\subsubsection{Edge Pruning in Low-Signal Regions}

\paragraph{Overlap count:} For each crown $c_{i,j}$, we compute the number of overlapping partners in $\mathcal{C}_{i+1}$:
\[
n_{\text{overlap}}^{\text{next}}(c_{i,j}) = \left|\left\{c_{i+1,k} : \max(\text{overlap}_{\text{prev}}, \text{overlap}_{\text{curr}}) \geq \tau_{\text{overlap}}\right\}\right|
\]
Similarly, for each $c_{i+1,k}$, we count overlapping partners in $\mathcal{C}_i$:
\[
n_{\text{overlap}}^{\text{prev}}(c_{i+1,k}) = \left|\left\{c_{i,j} : \max(\text{overlap}_{\text{prev}}, \text{overlap}_{\text{curr}}) \geq \tau_{\text{overlap}}\right\}\right|
\]

\paragraph{Low-signal pruning:} In regions with very low overlap (e.g., edge of orthomosaic, areas with poor alignment), many candidates may pass the spatial filter but have $\text{IoU} \approx 0$. To prevent spurious matches, we apply an additional pruning rule:
\begin{itemize}
    \item If $n_{\text{overlap}}^{\text{next}}(c_{i,j}) = 0$ and $\text{IoU}(c_{i,j}, c_{i+1,k}) < 0.10$, reject the candidate.
    \item If $n_{\text{overlap}}^{\text{prev}}(c_{i+1,k}) = 0$ and $\text{IoU}(c_{i,j}, c_{i+1,k}) < 0.10$, reject the candidate.
\end{itemize}
This ensures that zero-overlap matches are only retained if both crowns have no better alternatives.

%--------------------------------------------------------------------
\subsection{Edge Case Classification}
\label{sec:case_classification}

After computing features and applying early filters, each candidate edge is assigned to one of four mutually exclusive cases: \texttt{one\_to\_one}, \texttt{containment}, \texttt{nearby}, or \texttt{none}. The classification is performed using a decision tree based on overlap counts, IoU, overlap ratios, and containment flags.

%--------------------------------------------------------------------
\subsubsection{One-to-One Matches}

A candidate $(c_{i,j}, c_{i+1,k})$ is classified as \texttt{one\_to\_one} if:
\begin{enumerate}
    \item $n_{\text{overlap}}^{\text{next}}(c_{i,j}) = 1$ and $n_{\text{overlap}}^{\text{prev}}(c_{i+1,k}) = 1$ (unique overlap in both directions).
    \item $\text{overlap}_{\text{prev}}(c_{i,j}, c_{i+1,k}) \geq \tau_{\text{ovlp}}^{\text{1-1}}$ (high overlap from previous crown).
    \item $\text{overlap}_{\text{curr}}(c_{i,j}, c_{i+1,k}) \geq \tau_{\text{ovlp}}^{\text{1-1}}$ (high overlap to current crown).
    \item $\text{IoU}(c_{i,j}, c_{i+1,k}) \geq \tau_{\text{IoU}}^{\text{1-1}}$ (strong intersection over union).
\end{enumerate}

\paragraph{Parameter values:} For the strict preset:
\[
\tau_{\text{ovlp}}^{\text{1-1}} = 0.72, \quad \tau_{\text{IoU}}^{\text{1-1}} = 0.40
\]
For the ultra-relaxed preset:
\[
\tau_{\text{ovlp}}^{\text{1-1}} = \max(0.30, \tau_{\text{overlap}}), \quad \tau_{\text{IoU}}^{\text{1-1}} = \max(0.10, \tau_{\text{overlap}}/2)
\]

\paragraph{Interpretation:} \texttt{one\_to\_one} edges represent high-confidence matches where both crowns exhibit strong spatial coincidence and are each other's unique overlapping partner. These are the most reliable edges and are prioritized in selection.

\paragraph{Scoring:} For \texttt{one\_to\_one} edges, the final score is computed as:
\[
s^{\text{1-1}}(c_{i,j}, c_{i+1,k}) = \sum_{\beta} w_\beta^{\text{1-1}} \cdot f_\beta(c_{i,j}, c_{i+1,k})
\]
where $f_\beta \in \{s_{\text{base}}, \text{IoU}, \text{overlap}_{\text{prev}}, \text{overlap}_{\text{curr}}, f_{\text{centroid}}\}$ and typical weights are:
\[
w_{\text{base}}^{\text{1-1}} = 0.55, \quad w_{\text{IoU}}^{\text{1-1}} = 0.20, \quad w_{\text{ovlp,prev}}^{\text{1-1}} = 0.10, \quad w_{\text{ovlp,curr}}^{\text{1-1}} = 0.10, \quad w_{\text{centroid}}^{\text{1-1}} = 0.05
\]

%--------------------------------------------------------------------
\subsubsection{Splits and Merges}

While our current pipeline does not explicitly model split/merge edges as a separate case (to maintain one-to-one cardinality constraints), we can identify potential splits and merges post-hoc by examining graph topology:

\paragraph{Split pattern:} A crown $c_{i,j}$ is involved in a split if:
\[
|\{c_{i+1,k} : (c_{i,j}, c_{i+1,k}) \in E\}| > 1
\]
i.e., it has multiple outgoing edges. This occurs when relaxed cardinality constraints allow $\max_{\text{edges\_per\_prev}} > 1$.

\paragraph{Merge pattern:} A crown $c_{i+1,k}$ is involved in a merge if:
\[
|\{c_{i,j} : (c_{i,j}, c_{i+1,k}) \in E\}| > 1
\]
i.e., it has multiple incoming edges.

\paragraph{Handling in strict preset:} The strict preset enforces $\max_{\text{edges\_per\_prev}} = 1$ and $\max_{\text{edges\_per\_curr}} = 1$, eliminating split/merge edges. In ambiguous cases, only the highest-scoring edge is retained.

\paragraph{Handling in ultra-relaxed preset:} The ultra-relaxed preset permits up to 3 outgoing and 3 incoming edges per node, allowing splits/merges to be represented but subject to scoring thresholds to avoid excessive branching.

%--------------------------------------------------------------------
\subsubsection{Containment and Partial Overlaps}

A candidate $(c_{i,j}, c_{i+1,k})$ is classified as \texttt{containment} if:
\begin{enumerate}
    \item \textbf{Geometric containment}: $\text{prev\_contains\_curr} = \text{True}$ or $\text{curr\_contains\_prev} = \text{True}$.
    \item \textbf{High overlap}: $\text{overlap}_{\text{prev}} \geq \tau_{\text{ovlp}}^{\text{cont}}$ and $\text{overlap}_{\text{curr}} \geq \tau_{\text{ovlp}}^{\text{cont}}$.
\end{enumerate}

\paragraph{Parameter values:} For both presets:
\[
\tau_{\text{ovlp}}^{\text{cont}} = 0.82
\]
This high threshold ensures that containment edges are only created when one crown is nearly fully subsumed by the other (or vice versa), with minimal non-overlapping area.

\paragraph{Interpretation:} Containment edges arise when segmentation boundaries shift significantly between orthomosaics, but the spatial coincidence remains strong. Common causes include:
\begin{itemize}
    \item Refined segmentation in $\mathcal{O}_{i+1}$ that better delineates crown edges.
    \item Real canopy growth where a crown expands to encompass its previous extent.
    \item Segmentation artifacts where a large crown is incorrectly subdivided or merged.
\end{itemize}

\paragraph{Scoring:} Containment scores emphasize overlap ratios:
\[
s^{\text{cont}}(c_{i,j}, c_{i+1,k}) = w_{\text{base}}^{\text{cont}} \cdot s_{\text{base}} + w_{\text{ovlp,prev}}^{\text{cont}} \cdot \text{overlap}_{\text{prev}} + w_{\text{ovlp,curr}}^{\text{cont}} \cdot \text{overlap}_{\text{curr}}
\]
with typical weights:
\[
w_{\text{base}}^{\text{cont}} = 0.30, \quad w_{\text{ovlp,prev}}^{\text{cont}} = 0.35, \quad w_{\text{ovlp,curr}}^{\text{cont}} = 0.35
\]

%--------------------------------------------------------------------
\subsubsection{Proximity-Only Matches and \texttt{none} Edges}

\paragraph{\texttt{nearby} case (ultra-relaxed only):} A candidate is classified as \texttt{nearby} if:
\begin{enumerate}
    \item It does not qualify for \texttt{one\_to\_one} or \texttt{containment}.
    \item $d_{\text{centroid}}(c_{i,j}, c_{i+1,k}) < d_{\text{near}} = 30$~m (tight proximity).
    \item $\max(\text{overlap}_{\text{prev}}, \text{overlap}_{\text{curr}}) \geq \tau_{\text{overlap}}$ (minimal overlap requirement, often 5\%).
\end{enumerate}

\paragraph{Motivation:} The \texttt{nearby} case acts as a fallback for crowns that are spatially co-located but have poor shape similarity or IoU due to:
\begin{itemize}
    \item Large shape changes (e.g., deciduous trees transitioning from full foliage to partial leaf-off).
    \item Detection inconsistencies (boundary shifts, different confidence thresholds).
    \item Residual alignment errors.
\end{itemize}
By prioritizing spatial proximity, we can maintain chains even when overlap degrades.

\paragraph{Scoring:} Proximity-only scores are dominated by spatial similarity:
\[
s^{\text{near}}(c_{i,j}, c_{i+1,k}) = w_{\text{base}}^{\text{near}} \cdot s_{\text{base}} + w_{\text{centroid}}^{\text{near}} \cdot f_{\text{centroid}}
\]
with typical weights:
\[
w_{\text{base}}^{\text{near}} = 0.70, \quad w_{\text{centroid}}^{\text{near}} = 0.30
\]
The base similarity itself uses spatial weight 0.85, so proximity accounts for $\sim 90\%$ of the final score.

\paragraph{\texttt{none} case:} All candidates not meeting the criteria for \texttt{one\_to\_one}, \texttt{containment}, or \texttt{nearby} are assigned case \texttt{none} and discarded. No edge is created in the graph.

%--------------------------------------------------------------------
\subsection{Tracking Graph Construction and Statistics}

After case classification and scoring, we select edges to add to the graph $G$ using a case-prioritized greedy algorithm. This section describes the selection procedure and the resulting graph statistics.

%--------------------------------------------------------------------
\subsubsection{IoU Similarity Matrices Between Orthomosaic Pairs}

For each consecutive pair $(\mathcal{O}_i, \mathcal{O}_{i+1})$, we construct an $m_i \times m_{i+1}$ IoU matrix $\mathbf{M}_{i,i+1}$ where:
\[
[\mathbf{M}_{i,i+1}]_{jk} = \text{IoU}(c_{i,j}, c_{i+1,k})
\]
This matrix is visualized as a heatmap (Figure~\ref{fig:iou_matrices}) to diagnose alignment quality and overlap patterns.

\paragraph{Ideal matrix structure:} For perfectly aligned orthomosaics with consistent detection, $\mathbf{M}_{i,i+1}$ would have a clear diagonal structure (assuming crowns maintain their spatial order). In practice:
\begin{itemize}
    \item \textbf{Sparse structure}: Most entries are 0 due to spatial separation.
    \item \textbf{Off-diagonal nonzeros}: Indicate alignment drift or crown position changes.
    \item \textbf{Multiple high values per row}: Suggest splits (one crown in $\mathcal{O}_i$ overlaps multiple in $\mathcal{O}_{i+1}$).
    \item \textbf{Multiple high values per column}: Suggest merges.
\end{itemize}

\paragraph{Aggregate statistics:} For the SIT dataset (5 orthomosaics, ultra-relaxed preset), the IoU matrices exhibit:
\begin{itemize}
    \item Mean IoU (over nonzero entries): 0.42--0.58 depending on pair.
    \item Fraction of crown pairs with IoU $\geq 0.5$: 35--50\%.
    \item Median centroid distance for matched pairs: 8--25~m.
\end{itemize}

%--------------------------------------------------------------------
\subsubsection{Graph Degree Distributions and Connectivity Patterns}

\paragraph{Out-degree distribution:} The out-degree $d_{\text{out}}(v)$ of node $v = (i,j)$ is the number of outgoing edges:
\[
d_{\text{out}}(v) = |\{(v, u) \in E\}|
\]
For $v$ in orthomosaic $\mathcal{O}_i$ with $i < N$, we expect:
\begin{itemize}
    \item $d_{\text{out}}(v) = 0$: Crown $c_{i,j}$ has no match in $\mathcal{O}_{i+1}$ (disappeared or unmatched).
    \item $d_{\text{out}}(v) = 1$: Unique forward match (ideal for one-to-one tracking).
    \item $d_{\text{out}}(v) > 1$: Split or ambiguous match (allowed in ultra-relaxed, forbidden in strict).
\end{itemize}

\paragraph{In-degree distribution:} The in-degree $d_{\text{in}}(v)$ of node $v = (i+1,k)$ is the number of incoming edges:
\[
d_{\text{in}}(v) = |\{(u, v) \in E\}|
\]
Analogously:
\begin{itemize}
    \item $d_{\text{in}}(v) = 0$: Crown $c_{i+1,k}$ is newly appeared (no predecessor in $\mathcal{O}_i$).
    \item $d_{\text{in}}(v) = 1$: Unique backward match.
    \item $d_{\text{in}}(v) > 1$: Merge or ambiguous.
\end{itemize}

\paragraph{Empirical distributions (ultra-relaxed, 10Nov):}
\begin{center}
\begin{tabular}{lcccc}
\toprule
Degree & $d=0$ & $d=1$ & $d=2$ & $d \geq 3$ \\
\midrule
Out-degree & 114 (27.9\%) & 272 (66.7\%) & 18 (4.4\%) & 4 (1.0\%) \\
In-degree & 114 (27.9\%) & 276 (67.6\%) & 16 (3.9\%) & 2 (0.5\%) \\
\bottomrule
\end{tabular}
\end{center}
Most nodes have exactly one outgoing or incoming edge (one-to-one tracking). The $d=0$ nodes correspond to singletons (unmatched crowns). Higher-degree nodes are rare but represent potential split/merge events.

\paragraph{Average degrees:}
\[
\bar{d}_{\text{out}} = \frac{|E|}{|V|} = \frac{201}{408} \approx 0.49, \quad
\bar{d}_{\text{in}} = \frac{|E|}{|V|} \approx 0.49
\]
The low average degree reflects the sparse matching: only $\sim 63\%$ of crowns are matched to a successor/predecessor.

%--------------------------------------------------------------------
\subsubsection{Node and Edge Counts}

\paragraph{Total nodes:} $|V| = \sum_{i=1}^N |\mathcal{C}_i|$. For the SIT dataset:
\[
|V| = 82 + 78 + 78 + 78 + 80 = 408
\]

\paragraph{Total edges:} $|E| = 201$ (ultra-relaxed) vs.\ $|E| = 2$ (strict). The edge count reflects the trade-off between precision and recall:
\begin{itemize}
    \item Strict preset: Very few edges, high confidence, many broken chains.
    \item Ultra-relaxed preset: More edges, lower per-edge confidence, longer chains.
\end{itemize}

\paragraph{Edges by case (ultra-relaxed):}
\begin{center}
\begin{tabular}{lcc}
\toprule
Case & Candidates & Selected \\
\midrule
\texttt{one\_to\_one} & 122 & 122 (100\%) \\
\texttt{containment} & 8 & 5 (62.5\%) \\
\texttt{nearby} & 162 & 74 (45.7\%) \\
\bottomrule
\end{tabular}
\end{center}
All \texttt{one\_to\_one} candidates are selected (highest confidence). \texttt{nearby} edges have a 45.7\% selection rate due to competition and score thresholds.

\paragraph{Match rates by consecutive pair:}
\[
\text{match\_rate}(\mathcal{O}_i \to \mathcal{O}_{i+1}) = \frac{\text{\# edges from $\mathcal{O}_i$ to $\mathcal{O}_{i+1}$}}{\min(|\mathcal{C}_i|, |\mathcal{C}_{i+1}|)}
\]
For the SIT dataset (ultra-relaxed):
\begin{center}
\begin{tabular}{lcc}
\toprule
OM Pair & Matches & Match Rate \\
\midrule
$\mathcal{O}_1 \to \mathcal{O}_2$ & 57 / 82 & 0.695 \\
$\mathcal{O}_2 \to \mathcal{O}_3$ & 53 / 78 & 0.679 \\
$\mathcal{O}_3 \to \mathcal{O}_4$ & 50 / 78 & 0.641 \\
$\mathcal{O}_4 \to \mathcal{O}_5$ & 41 / 80 & 0.512 \\
\midrule
Overall & 201 / 318 & 0.632 \\
\bottomrule
\end{tabular}
\end{center}
Match rates degrade with increasing orthomosaic index, likely due to cumulative alignment drift and seasonal phenological changes.

%--------------------------------------------------------------------
\subsubsection{Chain Length Distribution}

\paragraph{Chain extraction:} We extract all maximal directed paths in $G$ by:
\begin{enumerate}
    \item Identifying source nodes: $S = \{v \in V : d_{\text{in}}(v) = 0\}$.
    \item For each $v \in S$, perform a greedy forward traversal:
    \begin{itemize}
        \item Start at $v$, set $\mathcal{P} = [v]$.
        \item While $d_{\text{out}}(v) > 0$:
        \begin{itemize}
            \item If $d_{\text{out}}(v) = 1$, append the unique successor to $\mathcal{P}$.
            \item If $d_{\text{out}}(v) > 1$, append the highest-scoring successor to $\mathcal{P}$ (greedy).
        \end{itemize}
        \item Stop when reaching a sink ($d_{\text{out}}(v) = 0$).
    \end{itemize}
    \item Output $\mathcal{P}$ as a chain.
\end{enumerate}

\paragraph{Chain length:} The length $\ell$ of chain $\mathcal{P}$ is the number of nodes. For $N$ orthomosaics, $\ell \in \{1, 2, \ldots, N\}$.

\paragraph{Distribution (ultra-relaxed, 10Nov):}
\begin{center}
\begin{tabular}{lcc}
\toprule
Chain Length $\ell$ & Count & Percentage \\
\midrule
1 & 114 & 27.9\% \\
2 & 39 & 9.6\% \\
3 & 21 & 5.1\% \\
4 & 12 & 2.9\% \\
5 (full) & 21 & 5.1\% \\
\bottomrule
\end{tabular}
\end{center}

\paragraph{Summary statistics:}
\begin{align*}
\text{Total chains} &= 207 \\
\text{Average chain length} &= 1.97 \\
\text{Median chain length} &= 1.0 \\
\text{Max chain length} &= 5 \ (\text{full coverage across all OMs})
\end{align*}

\paragraph{Interpretation:}
\begin{itemize}
    \item \textbf{Length-1 chains (singletons)}: 27.9\% of crowns are unmatched. These may be false positive detections, crowns at the edge of orthomosaics, or trees with unique temporal signatures (e.g., newly planted, removed).
    
    \item \textbf{Length-2 chains}: 9.6\% tracked for only two consecutive OMs. Could represent ephemeral detections or trees on the study area boundary.
    
    \item \textbf{Length-3 and length-4 chains}: 8\% tracked for 3--4 OMs. Possible deciduous trees with intermittent detection (leaf-off periods) or crowns with inconsistent segmentation.
    
    \item \textbf{Length-5 chains (full)}: 5.1\% successfully tracked across all 5 orthomosaics. These are likely evergreen trees with stable canopy structure and high detection consistency. These chains form the gold standard for phenology analysis.
\end{itemize}

%--------------------------------------------------------------------
\subsubsection{Examples of Good and Bad Chains}

\paragraph{Example 1: High-quality full-length chain.}
Consider chain $\mathcal{P}_1 = [(1, 34), (2, 29), (3, 31), (4, 28), (5, 35)]$ with edge scores:
\begin{center}
\begin{tabular}{lcccc}
\toprule
Edge & Case & Similarity & IoU & Centroid Dist (m) \\
\midrule
$(1,34) \to (2,29)$ & \texttt{one\_to\_one} & 0.89 & 0.72 & 5.2 \\
$(2,29) \to (3,31)$ & \texttt{one\_to\_one} & 0.91 & 0.78 & 3.8 \\
$(3,31) \to (4,28)$ & \texttt{one\_to\_one} & 0.87 & 0.68 & 6.1 \\
$(4,28) \to (5,35)$ & \texttt{one\_to\_one} & 0.90 & 0.74 & 4.5 \\
\bottomrule
\end{tabular}
\end{center}
\textbf{Characteristics}: All edges are \texttt{one\_to\_one} with high similarity ($>0.87$) and strong IoU ($>0.68$). Centroid distances are small ($<7$~m), indicating minimal alignment drift. This chain likely represents a stable evergreen tree with consistent detection.

\paragraph{Example 2: Moderate-quality partial chain.}
Consider chain $\mathcal{P}_2 = [(1, 12), (2, 8), (3, 15)]$ with edge scores:
\begin{center}
\begin{tabular}{lcccc}
\toprule
Edge & Case & Similarity & IoU & Centroid Dist (m) \\
\midrule
$(1,12) \to (2,8)$ & \texttt{one\_to\_one} & 0.75 & 0.52 & 12.3 \\
$(2,8) \to (3,15)$ & \texttt{nearby} & 0.48 & 0.18 & 18.7 \\
\bottomrule
\end{tabular}
\end{center}
\textbf{Characteristics}: First edge is \texttt{one\_to\_one} but with moderate quality. Second edge is \texttt{nearby} with low IoU and larger centroid distance. The chain breaks after $\mathcal{O}_3$ (crown $(3,15)$ has no successor). This pattern suggests either:
\begin{itemize}
    \item Detection failure in $\mathcal{O}_4$ and $\mathcal{O}_5$ (deciduous tree with leaf-off).
    \item Large shape change between $\mathcal{O}_2$ and $\mathcal{O}_3$ breaking similarity.
    \item Alignment drift increasing centroid separation.
\end{itemize}

\paragraph{Example 3: Low-quality chain (potential false positive).}
Consider chain $\mathcal{P}_3 = [(2, 45), (3, 58)]$ with edge score:
\begin{center}
\begin{tabular}{lcccc}
\toprule
Edge & Case & Similarity & IoU & Centroid Dist (m) \\
\midrule
$(2,45) \to (3,58)$ & \texttt{nearby} & 0.32 & 0.08 & 25.4 \\
\bottomrule
\end{tabular}
\end{center}
\textbf{Characteristics}: Length-2 chain with a single \texttt{nearby} edge. Very low similarity and IoU, large centroid distance. This chain likely represents a spurious match between unrelated crowns or transient detection artifacts. Such chains would be filtered in a stricter preset or post-processing step.

\paragraph{Example 4: Branching chain (ultra-relaxed, split scenario).}
Consider node $(2, 15)$ with two outgoing edges:
\begin{center}
\begin{tabular}{lcccc}
\toprule
Edge & Case & Similarity & IoU & Centroid Dist (m) \\
\midrule
$(2,15) \to (3,22)$ & \texttt{one\_to\_one} & 0.68 & 0.45 & 9.1 \\
$(2,15) \to (3,23)$ & \texttt{containment} & 0.62 & 0.38 & 11.7 \\
\bottomrule
\end{tabular}
\end{center}
\textbf{Characteristics}: Crown $(2,15)$ branches into two successors in $\mathcal{O}_3$. Visual inspection (RGB patches) reveals that crown $(2,15)$ is a large canopy that over-segments into two distinct crowns $(3,22)$ and $(3,23)$ in $\mathcal{O}_3$ due to improved segmentation detail. The greedy chain extraction selects the \texttt{one\_to\_one} edge as the primary successor, but both edges are retained in the graph for analysis.

\paragraph{Summary:} High-quality chains (full-length, all \texttt{one\_to\_one} edges, high similarity) constitute $\sim 5\%$ of the total and represent the most reliable tracking. Moderate chains with \texttt{nearby} edges or partial length are more common ($\sim 20\%$) and require careful interpretation. Low-quality chains and singletons ($\sim 28\%$) may include false positives or transient detections.

%--------------------------------------------------------------------
\subsection{Discussion: Trade-offs and Parameter Sensitivity}

\paragraph{Precision vs.\ recall:} The strict preset (Section~\ref{sec:strict_preset}) achieves very high precision (nearly all edges are correct matches) but extremely low recall (only 0.4\% match rate). The ultra-relaxed preset (Section~\ref{sec:relaxed_preset}) achieves 63.2\% recall at the cost of lower precision (some \texttt{nearby} edges are likely false positives). For phenology analysis requiring long-term tracking, the ultra-relaxed preset is preferred despite lower per-edge confidence.

\paragraph{Threshold tuning:} Small changes in $\tau_{\text{overlap}}$ or $\tau_{\text{sim}}$ dramatically affect results. For example, increasing $\tau_{\text{overlap}}$ from 0.05 to 0.10 reduces edge count by $\sim 20\%$, while increasing $\tau_{\text{ovlp}}^{\text{1-1}}$ from 0.72 to 0.80 reduces \texttt{one\_to\_one} edges by $\sim 40\%$. This sensitivity underscores the need for application-specific parameter tuning guided by validation against ground truth (when available).

\paragraph{Alignment impact:} Spatial alignment (Section~3) improves match rates by 5--15\% on average, with the greatest benefit for mid-sequence pairs ($\mathcal{O}_2 \to \mathcal{O}_3$). First pairs ($\mathcal{O}_1 \to \mathcal{O}_2$) benefit less due to inherently better overlap. Late pairs ($\mathcal{O}_4 \to \mathcal{O}_5$) benefit modestly, as cumulative drift exceeds what simple translational alignment can correct.

\paragraph{Case priority:} Processing cases in order (\texttt{one\_to\_one}, \texttt{containment}, \texttt{nearby}) ensures high-confidence edges are selected first, preventing lower-quality matches from preempting better alternatives. Reversing this order (processing \texttt{nearby} first) would increase edge count but substantially degrade precision.

\paragraph{Computational cost:} Candidate generation is $O(|\mathcal{C}_i| \times |\mathcal{C}_{i+1}|)$ per pair, but spatial filtering reduces this to $O(k \cdot |\mathcal{C}_i|)$ where $k \ll |\mathcal{C}_{i+1}|$ is the average number of spatially proximate crowns. Edge selection is $O(|E_{\text{candidates}}| \log |E_{\text{candidates}}|)$ due to sorting by score. For the SIT dataset, total runtime is $<5$ minutes on a standard laptop (graph construction dominates).

\paragraph{Future enhancements:}
\begin{itemize}
    \item \textbf{Multi-date context}: Incorporate information from $\mathcal{O}_{i-1}$ and $\mathcal{O}_{i+2}$ to disambiguate matches at $\mathcal{O}_i \to \mathcal{O}_{i+1}$.
    \item \textbf{Learned features}: Train a neural network to embed crown features, replacing hand-crafted similarity metrics.
    \item \textbf{Virtual nodes}: Insert virtual nodes for crowns that disappear temporarily (e.g., deciduous trees in leaf-off), enabling chain continuation across gaps.
    \item \textbf{Global optimization}: Formulate tracking as a min-cost flow problem on $G$ to enforce global chain consistency~\cite{zhang2008global}.
\end{itemize}
